% The checks behind the paper's computations, each with the control that must fail.
% \input after the exponential's section of expansion-error-bounds.tex, whose labels it uses.

\section{The checks and their controls}\label{sec:checks}

Each computation behind this paper comes with a control: the same check on a wrong result, a wrong input or a weakened
algorithm, which must fail. Table~\ref{tab:checks} lists the checks and their controls.

\begin{table*}[tp]
\centering
\caption{The checks, and the control of each that must fail.}
\label{tab:checks}
\footnotesize
\begin{tabular}{@{}>{\raggedright\arraybackslash}p{0.105\textwidth}>{\raggedright\arraybackslash}p{0.545\textwidth}>{\raggedright\arraybackslash}p{0.3\textwidth}@{}}
\toprule
check & what is checked & control \\
\midrule
constants & Each constant, from its proved recursion (\texttt{constants.py}, \texttt{div\_consts.py},
\texttt{sqrt\_consts.py}, \texttt{acc\_bound.py}, \texttt{sloppy\_consts.py}, \texttt{trig\_consts.py},
\texttt{explog\_consts.py}). & The recomputation below. \\ \addlinespace
recomputation & A second set of programs (\texttt{review/}), written from the text of the appendix before its author
read those scripts and sharing none of their code, recomputes Table~\ref{tab:kappa}, $\kdiv$, the square root's
bound, and the constants of Theorems~\ref{thm:sloppy}, \ref{thm:trig}, \ref{thm:exp} and~\ref{thm:log} and of
Lemma~\ref{lem:lnb}, rounding up at 300 bits. It evaluates every bucket of Theorem~\ref{thm:trig} and the cases of
Theorems~\ref{thm:exp} and~\ref{thm:log} as their statements define them, and agrees with every digit quoted in this
paper: $3.163167393\,u^8$ and $6.418350521\,u^8$ for the sine, cosine and tangent, $3.660060455\,u^8$ for the
exponential, $9.925819\,u^8$ and $2^{-509.0880}$ for the logarithm, and $5.914423\,u$ for the binary64 logarithm. &
Each program has a control that must fail. \\ \addlinespace
bits at $n = 8$ & At $p = 53$ the emulator returns multidouble's bits on 10,000 inputs of each of ADD, ADDF, MULF,
MUL and DIV: operands with dense parts, parts at the top of the ladder, gaps and zero parts, and cancelling pairs for
ADD. Every error is within its bound, at most $0.0073$ of it. Every product, and every sum of operands whose leading
parts have one sign, is on the ladder. Dense random operands err by at most $0.0064\,u^8|Z|$ (\texttt{campaign.py}). &
An emulator with one error term dropped mismatches on 182 to 300 of 300 inputs. \\ \addlinespace
renormalization & At $p = 3$ and $n = 3$ the renormalization of every four terms of a family of 61 values, in every
order, 7,036,411 tuples, loses at most $\min_k B_k$, only after $n - 1$ emissions, and exactly $\min_k B_k$ on 40,600
of them. Every three terms of 81 values are renormalized exactly. A second pass with $\FastTwoSum$ throughout loses
more on 948,874 of the four-term tuples. & Half the bound fails on 281,138 tuples. \\ \addlinespace
small precisions & At $(p, n) = (5, 3)$ and $(7, 4)$ the five operations stay within Theorems~\ref{thm:t3}
and~\ref{thm:t6} on 1,166,592 and 352,000 pairs of operands whose parts are at the edge of the ladder, just below
$u^i$, far below, or zero, with both signs. & With a $\TwoSum$ whose error has the wrong sign, most pairs break the
bound. \\ \addlinespace
sharp inputs & The inputs of Proposition~\ref{prop:sharp} give its errors exactly for $n = 2$ to $8$, and multidouble
returns $\pm(1, u, \dots, u^7)$ on them at $n = 8$. & Half the bound fails. \\ \addlinespace
example & The example of Section~\ref{sec:example} runs on the same emulator (\texttt{example\_n3.py}). & Its control
must fail. \\ \addlinespace
square root & The emulator's square root matches multidouble's bit for bit on 2000 inputs with gaps, zero parts and
parts at the top of the ladder (\texttt{sqrt\_bits.py}). & Skipping its last step mismatches on every input, and
skipping the two-sum of the leading parts on 30. \\ \addlinespace
accumulator & The bound, in exact rational arithmetic, holds on the constructed inputs of
Proposition~\ref{prop:acc-low} at $p = 8$ to $14$ and at $p = 53$, and on random sums (\texttt{accumulator/}).
Multidouble returns the same eight parts as the emulator on the inputs of Proposition~\ref{prop:acc-low}. & A bound
with the powers of $K$ lowered by one, a bound linear in $K$, and an accumulator with one level's error dropped each
fire. \\ \addlinespace
exponential and logarithm & \texttt{explog\_consts.py} samples the convexity that the logarithm's chords use, on
every case at $k = 0$ and four at $k \ne 0$. The emulator's operations run EXP and LOG as the appendix states them
(\texttt{explog\_bits.py}) and give multidouble's \texttt{OD::exp}, \texttt{OD::ln}, \texttt{DD::exp} and
\texttt{DD::ln} bit for bit on 3996 arguments, random, near 0, near multiples of $\frac12\ln 2$, near 1, subnormal,
with parts at the top of the ladder, and for the logarithm within 3 ulps of the largest finite value; their results are
within the bounds against Python's correctly rounded exp and ln at 1300 digits. The same script hooks every product,
fused multiply-add and quotient of the steps to find which leave the normal range, and how far the closest clears it
(Appendix, \emph{Which runs meet} (A2$'$)). & A concave function, for the convexity. Halving the Taylor degree, or
moving the logarithm's start by $2^{-20}$, mismatches and breaks the bounds. At full Taylor degree the hook must find
products that leave the range. \\ \addlinespace
inputs of Table~\ref{tab:inputs} & Checked from scratch (\texttt{table1\_inputs.py}): the emulator's results on
them, the exact values enclosed by code that shares nothing with the algorithms (an integer square root, Taylor series
of the sine and cosine in fixed point, and Python's correctly rounded decimal exp and ln at 1300 digits), the ratios of
the errors to $u^8|Z|$, and multidouble's bits on the same inputs, which agree with the emulator's. & Rounding ties
away from zero, which breaks the chains of ties, and moving each input's leading part up by one ulp: each drops
every ratio far below its stated value. \\ \addlinespace
multidouble's tests & The bound of Theorem~\ref{thm:main} against exact arithmetic on every random input that meets
the hypotheses, with the inputs of Proposition~\ref{prop:worst} and Table~\ref{tab:inputs} pinned; the bounds of the
exponential and the logarithm against a reference accurate to $2^{-1600}$. & The reference is checked against planted
truth. \\ \addlinespace
Lean & The statements of the appendices are proved in Lean~4, and its kernel evaluates their constants
(Appendix~\ref{sec:lean}). Tables~\ref{tab:lean} and~\ref{tab:lean2} are checked against the sources: every file
at the top of \texttt{ExpansionBounds/} is in a row, every folder under it is named in that appendix, and every
theorem named is declared in its row's files (\texttt{lean\_refs.py}). & The build reports a theorem that rests on a
planted axiom, and the kernel proves each check false on a wrong number or state (\texttt{Controls/}). The tables'
check reports a theorem and a file that do not exist, a folder that the appendix does not name, and the files of
Table~\ref{tab:lean2} when it is left out. \\
\bottomrule
\end{tabular}
\end{table*}
